\documentclass[11pt]{article}
\usepackage[margin=1in]{geometry}
\usepackage{amsmath,amssymb,amsthm}
\usepackage[colorlinks=true,urlcolor=blue]{hyperref}
\newtheorem{theorem}{Theorem}
\newtheorem{proposition}[theorem]{Proposition}
\newtheorem{remark}[theorem]{Remark}
\newcommand{\one}{\boldsymbol 1}
\newcommand{\R}{\mathbb R}
\newcommand{\eps}{\varepsilon}
\title{One edge: two exact bisection classes,\\but arbitrarily many orthogonal Fiedler bisections}
\author{Research-session theorem and obstruction; priority not established}
\date{November 15, 2026}
\begin{document}
\maketitle
\begin{abstract}
Increasing one edge conductance produces an exact minimum-bisection
value with at most two affine pieces: one fixed class of cuts crosses
the edge, and the other does not. In contrast, for every $k$, one
existing-edge doubling in a positive-weight complete graph can pass
through $k$ pairwise orthogonal balanced Fiedler sign cuts. The Fiedler
eigenvalue is simple throughout, and at the selected parameters every
coordinate has magnitude comparable to $n^{-1/2}$. Baseline edge weights
lie in $[1,3]$, all weights throughout the doubling lie in $[1,6]$, and
the relative Laplacian perturbation is at most $6/(2n-1)$.
Consequently, at any prescribed positive upper tolerance on relative
operator error, there is no finite graph-size-independent dimension
bound describing all Fiedler bisections along one edit.
This is an obstruction to a proposed gap-free route, not an improved
partitioning algorithm or a claim of a field-scale breakthrough.
\end{abstract}

\section*{Scope and exact statement}
All graphs are undirected, have unit vertex masses, and have positive
edge conductances. Bisection means exactly $n/2$ vertices on each side,
and the objective is total crossing conductance. The Laplacian is the
unnormalized Laplacian. A Fiedler vector is a unit eigenvector for its
second-smallest eigenvalue. A balanced sign cut is also its unambiguous
median bisection whenever exactly half of its coordinates have each sign.
No statement below substitutes a cardinality constraint for a
degree-volume constraint in normalized conductance.

\begin{proposition}[Exact one-edge cut envelope]\label{prop:envelope}
Fix a graph, an edge $e$, and any fixed finite family $\mathcal F$ of
admissible unoriented cuts. Increase only the conductance of $e$ by
$t\ge0$. Write $\mathcal F_i$ for cuts of edge-crossing status
$i\in\{0,1\}$, and let $a_i$ be their minimum baseline cost
(with $a_i=+\infty$ for an empty class). Then the optimum cost is
\[
                        \min\{a_0,a_1+t\}.
\]
The full optimizer set is the baseline minimizing set in the winning
class, or the union of both minimizing sets at a tie. In particular two
fixed cut representatives suffice for every $t$, and there is at most
one transition between the classes. For $r$ edited edges, at most
$2^r$ fixed representatives suffice for every vector of nonnegative
conductance increases.
\end{proposition}
\begin{proof}
Every cut in $\mathcal F_i$ has updated cost equal to its baseline
cost plus $it$. For several edges the increment is the scalar product
of the parameter vector with the cut's $\{0,1\}^r$ crossing signature.
Minimize separately over signatures and then over their at most $2^r$
classes. This is an elementary finite affine-envelope identity; its
algorithmic use requires computing the constrained optima $a_i$.
\end{proof}

\begin{theorem}[Unbounded spectral trajectory during one doubling]
\label{thm:trajectory}
For every integer $k\ge1$, every $0<\eta<1/10$, and every sufficiently
large power of two $n$ with $n/2-1\ge k$, there is a complete weighted
graph $G$ on $n$ vertices and an existing edge $e=uv$ with baseline
conductance $w_e$ such that the following hold. Put
$b=e_u-e_v$, $A=L(G)$, and $A_t=A+tbb^T$ for $0\le t\le w_e$.
\begin{enumerate}
\item All baseline conductances lie in $[1,3]$, and all conductances
throughout the doubling lie in $[1,6]$.
\item The Fiedler eigenvalue of $A_t$ is simple for every $t\ge0$.
There are $0<t_1<\cdots<t_k<w_e$, and orthonormal balanced sign
vectors $q_1,\ldots,q_k\in\{\pm n^{-1/2}\}^n$, such that a unit
Fiedler vector $f_j$ of $A_{t_j}$ satisfies
\[
 \|f_j-q_j\|_2\le2\eta,\qquad
 \operatorname{sign}(f_j)=\sqrt n\,q_j,
\]
and every coordinate satisfies
\[
 \frac{1-\eta}{(1+\eta)\sqrt n}
 \le |f_j(v)|\le
 \frac{1+\eta}{(1-\eta)\sqrt n}.
\]
Thus every pair of selected spectral bisections differs on exactly
$n/2$ vertices, even after choosing their orientations optimally.
\item On $\one^\perp$, uniformly throughout the doubling,
\[
 A\preceq A_t\preceq(1+\delta_n)A,
 \qquad \delta_n=\frac6{2n-1}.
\]
For every $0<\alpha\le1$, therefore,
\[
 0\preceq A^{-\alpha/2}(A_t^\alpha-A^\alpha)A^{-\alpha/2}
 \preceq\bigl((1+\delta_n)^\alpha-1\bigr)I.
\]
\end{enumerate}
The same $A_t$ eigenvectors are used by every positive power $A_t^\alpha$;
fractional powers do not change the spectral cut sequence.
\end{theorem}

\begin{proof}
\textit{Flat cut vectors and one rotation.}
Identify the vertices with $\{0,1\}^m$, $n=2^m$, and take
$u=(0,\ldots,0)$, $v=(1,0,\ldots,0)$. The $n/2-1$ nonconstant
Walsh characters not depending on the first coordinate give orthonormal
balanced sign vectors with $q_j(u)=q_j(v)=n^{-1/2}$. Select $k$ of
them. Put $e_0=b/\sqrt2$, which is orthogonal to these vectors and
to $\one$.

Set $c=1/(10k)$. Choose $0<r<1$ so small that
\begin{equation}\label{eq:r}
                 \frac{4\sqrt r}{1-\sqrt r}<\frac\eta4,
\end{equation}
and, subsequently, choose $d_1>0$ sufficiently small. Put
$d_j=d_1r^{j-1}$, with $d_1<1/100$. Define
\[
 \beta_j=\sqrt{cd_j/2},\quad
 \beta=\Bigl(\sum_j\beta_j^2\Bigr)^{1/2},\quad
 z=\beta^{-1}\sum_j\beta_jq_j,\quad
 \gamma=\sqrt{1-\beta^2}.
\]
Rotate only the plane spanned by $e_0,z$ as follows:
\[
 v_0=\gamma e_0-\beta z,\qquad
 v_j=q_j+(\gamma-1)\frac{\beta_j}{\beta}z+\beta_je_0.
\]
The vectors $v_0,v_1,\ldots,v_k$ are orthonormal in $\one^\perp$,
and
\begin{equation}\label{eq:couplings}
 b^Tv_0=\sqrt{2-c\sum_\ell d_\ell}=:b_0,
 \qquad b^Tv_j=\sqrt{cd_j}=:b_j.
\end{equation}
The vector $b$ has no component outside their span. With $k,n,r$
fixed, $v_j\to q_j$ and $v_0\to e_0$ as $d_1\downarrow0$.

\textit{Graph realization.}
Let $P=I-\one\one^T/n$, and let $P_R$ project onto the orthogonal
complement of $v_0,\ldots,v_k$ within $\one^\perp$. Define
\[
 K=2n,\qquad A=KP-v_0v_0^T+\sum_{j=1}^k d_jv_jv_j^T+P_R.
\]
The perturbation $M=A-KP$ has norm at most one and annihilates
$\one$. Thus $A$ has row sums zero, and its off-diagonal entries
satisfy $-3\le A_{xy}\le-1$. It is exactly the Laplacian of a
complete graph with conductances $-A_{xy}\in[1,3]$. On $\one^\perp$
its eigenvalues are $K-1$ at $v_0$, $K+d_j$ at $v_j$, and $K+1$
on the remaining subspace.

\textit{The secular equation.}
For $K-1<\mu<K+d_k$, put
\[
 m(\mu)=b^T(A-\mu I)^{-1}b
 =\frac{b_0^2}{K-1-\mu}+\sum_{\ell=1}^k
                         \frac{cd_\ell}{K+d_\ell-\mu}.
\]
Here the inverse is on $\one^\perp$. It is strictly increasing from
$-\infty$ to $+\infty$, since its derivative is a sum of positive
squares. For every $t>0$, the equation $1+tm(\mu)=0$ has a unique
solution in this interval, and the associated eigenvector is
$(A-\mu I)^{-1}b$. Rank-one interlacing makes this eigenvalue the
smallest one on $\one^\perp$. It is simple, since the secular root
is simple and all active poles are simple with nonzero couplings.
At $t=0$ the lowest eigenvalue $K-1$ is also simple.

Take $\mu_j=K-d_j$ and set $t_j=-1/m(\mu_j)$. At these parameters,
\[
 m(\mu_j)=-\frac{2-c\sum_\ell d_\ell}{1-d_j}
             +\sum_{\ell=1}^k\frac{cd_\ell}{d_\ell+d_j}<-1.
\]
Indeed $c\sum d_\ell<1/1000$ whereas the positive sum is at most
$ck=1/10$. Thus $0<t_j<1\le w_e$, and strict monotonicity of $m$
gives $t_1<\cdots<t_k$.

\textit{Dominance and rounding.}
Write the unnormalized eigenvector at $\mu_j$ as
\[
 y_j=-\frac{b_0}{1-d_j}v_0+
                  \sum_{\ell=1}^k\frac{\sqrt{cd_\ell}}{d_\ell+d_j}v_\ell.
\]
The coefficient of $v_j$ is $a_j=\sqrt c/(2\sqrt{d_j})$. For
$\ell\ne j$, the ratio of its positive coefficient to $a_j$ is
\[
 \frac{2\sqrt{d_\ell d_j}}{d_\ell+d_j}
             \le2r^{|\ell-j|/2}.
\]
Their sum is less than $\eta/4$ by~\eqref{eq:r}. The $v_0$
coefficient divided by $a_j$ tends to zero as $d_1\downarrow0$,
uniformly over the finite set of indices. Also
$\max_j\|v_j-q_j\|_\infty\to0$. Choose $d_1$ small enough that
the $v_0$ contribution and the sum of rotation errors each have
supremum norm less than $\eta/(4\sqrt n)$. The other $q_\ell$
contributions have supremum norm less than $\eta/(4\sqrt n)$.
Consequently
\[
                   \|y_j/a_j-q_j\|_\infty<\eta/\sqrt n.
\]
This gives the exact signs, $1-\eta\le\|y_j/a_j\|_2\le1+\eta$,
the displayed coordinate bounds after normalization, and
$\|y_j/\|y_j\|_2-q_j\|_2\le2\eta$.

\textit{Relative error.}
Since $A\succeq(2n-1)P$, $\|b\|^2=2$, and $t\le w_e\le3$,
\[
 tbb^T\preceq6P\preceq\frac6{2n-1}A.
\]
The fractional assertion follows from the Loewner--Heinz inequality
on $\one^\perp$. Orthogonality of balanced sign vectors makes the
associated cuts differ on $n/2$ vertices. Finally a positive power
preserves every Laplacian eigenspace and eigenvalue ordering.
\end{proof}

\begin{proposition}[Dimension lower bound]\label{prop:dimension}
For the selected cuts in Theorem~\ref{thm:trajectory}, any linear
subspace $V\subset\R^n$ satisfying
$\operatorname{dist}(q_j,V)^2\le\eps$ for every $j$ obeys
\[
                         \dim V\ge(1-\eps)k.
\]
If instead it approximates each selected unit Fiedler vector with
distance at most $\rho$, then
$\dim V\ge k(1-(\rho+2\eta)^2)$ whenever $\rho+2\eta<1$.
\end{proposition}
\begin{proof}
For an orthonormal list $q_j$ and an orthogonal projector $P_V$,
$\sum_{j=1}^k\|P_Vq_j\|^2\le\operatorname{tr}P_V=\dim V$.
The second assertion follows from
$\operatorname{dist}(q_j,V)\le\rho+2\eta$.
\end{proof}

\section*{What this settles, and what remains open}
The cut-envelope proposition and the trajectory theorem distinguish
the exact objective for cardinality bisections from its spectral
relaxation. During one edit, at most two fixed exact cut classes are
needed, while the spectral heuristic can visit an unbounded number
of orthogonal, unambiguously rounded bisections. A low-rank relative
approximation of a fractional power is a quadratic-form statement.
It cannot alone bound the dimension of the Fiedler-cut trajectory,
even if its error tends to zero.

More precisely, for every fixed $\eps>0$ and every integer $k$, one can
choose $n$ so that $\delta_n\le\eps$ and the trajectory still contains
$k$ orthogonal cut vectors. Thus there is no finite uniform dimension
bound for all graphs whose total relative perturbation is \emph{at most}
$\eps$. We do not claim to exclude every possible nonmonotone function
of the \emph{exact} perturbation norm; such a function could diverge
as that norm approaches zero.

The construction uses dense weighted graphs and arbitrarily small
spectral gaps. It does not prove the same many-cut phenomenon for
simple unweighted bounded-degree graphs. It does not make $k$
different exact optimum bisections win during one edge increase;
Proposition~\ref{prop:envelope} forbids that. It also does not show
that any of the selected cuts is a sparsest cut for a normalized
conductance objective, or refute useful gap-dependent perturbation
theorems. All bisections in the dense realization are within a
$1+O(1/n)$ objective factor of one another, so this is primarily an
embedding-complexity obstruction.

An affirmative substantial next target would be a gap-free,
computable certificate approximating the \emph{entire robust balanced
cut family}, whose complexity is controlled by a combinatorial
margin or cut-signature parameter rather than by a fractional
response rank. The elementary envelope does not supply an efficient
method to obtain its two constrained optima.

\section*{Primary literature and priority boundary}
Severe failures of spectral bisection, including failures of multiple
low eigenvectors, are classical; see Guattery and Miller,
\emph{On the Quality of Spectral Separators}, SIAM J. Matrix Anal.
Appl. 19 (1998),
\url{https://www.cs.cmu.edu/~glmiller/Publications/GuatteryMiller98.pdf}.
Theorem~\ref{thm:trajectory} is narrower than a new general failure
of spectral partitioning: its additional object is one fixed positive
rank-one edge trajectory, with many orthogonal flat rounded outputs.

Peng and Yoshida study average sensitivity of spectral clustering
under random edge deletions, with spectral-structure hypotheses:
\url{https://arxiv.org/abs/2006.04094}.
Their framework does not assert adversarial gap-free label stability.
Noschese and Reichel study structured single-edge sensitivity
of Perron and Fiedler eigenpairs (2025):
\url{https://onlinelibrary.wiley.com/doi/10.1002/nla.70041}.
The secular equation and rank-one interlacing used here are classical
tools, not a new perturbation calculus.

The exact rank-one eigenvector resolvent formula also appears as
Equation~(13) in the rank-one diagonalization paper of \v{S}oda and
Kempf (2025), arXiv:2206.09927v2:
\url{https://arxiv.org/html/2206.09927v2}.
Gabbassov and Kempf, \emph{Adiabatic Dynamics of Entanglement},
Quantum Science and Technology 10, 045054 (2025), describe nearly
orthogonal eigenvector exchanges and their cascades:
\url{https://arxiv.org/abs/2407.18301}.
Theorem~1 and Corollary~1.1 of the follow-up by Gabbassov, Singh and
Kempf, arXiv:2607.14227 (2026), state this cascade explicitly:
\url{https://arxiv.org/html/2607.14227v1}.
Their cascade transports a direction across successive adjacent
\emph{eigenvalue indices}. Our construction keeps the \emph{lowest
nontrivial index fixed} and places many different approximate directions
within its first interlacing interval. The graph-edge realization,
balanced flat rounding, and vanishing total relative norm are additional
features here. None of these observations makes the underlying
rank-one resolvent mechanism new.

d'Orsi, Jones, Ruotolo, Vadhan and Zhang (2025) formalize sparsest-cut
solution dimension and spectral enumeration plus cut improvement:
\url{https://arxiv.org/abs/2412.17115}.
Their low-energy projection and net arguments already address
near-optimal cut spaces under appropriate spectral or graph structure.
The elementary tail estimate for cut vectors and a repackaging of
their solution-dimension framework would not be a breakthrough.
No inspected primary theorem was found stating precisely the
many-orthogonal-output one-edge trajectory proved above. This limited
search is not a historical priority certification.
\end{document}
